Esta tesis presenta, implementa y compara cuatro maneras de representar el contexto de una melodía. La pregunta que guía el trabajo es sencilla: cuando varios modelos reciben la misma sucesión de notas, ¿la complejidad adicional permite anticipar mejor la nota siguiente y compensa el costo necesario para obtener esa mejora?

Se comparan cuatro alternativas: modelos de conteos, modelos con un número limitado o abierto de estados y un modelo neuronal capaz de aprovechar más contexto. Todas comparten la misma representación, y los arreglos de cada obra permanecen en una misma partición. De los $3000$ archivos seleccionados de PDMX, una colección pública de partituras digitales, quedaron $2933$ piezas después de preparar los datos.

Con todos los datos de entrenamiento, el modelo neuronal asignó la mayor probabilidad a las notas de prueba, y la comparación por obra confirmó su ventaja. El orden se mantuvo ante cambios en las ventanas, la capacidad del modelo finito y la partición de las obras. Eliminar el solapamiento redujo la ventaja neuronal; ampliar la capacidad finita mejoró su resultado, sin alcanzar un punto de saturación.

En la escala estudiada, el modelo neuronal ofreció una ventaja predictiva clara, pero no fue superior en todos los sentidos. Su ajuste tomó menos tiempo que dos alternativas probabilísticas y más que el modelo de conteos; además, el hardware especializado elevó el costo estimado bajo las tarifas documentadas. El prototipo se ejecutó de principio a fin en el equipo disponible. Como no hubo evaluación con usuarios ni anotaciones estructurales, los resultados sólo comparan predicción y costo.
